\documentclass[UTF8,fontset=fandol,11pt,a4paper]{ctexart}
\usepackage[a4paper,margin=1.72cm,headheight=14pt]{geometry}
\usepackage{amsmath,amssymb,mathtools}
\usepackage{booktabs,tabularx,longtable,array,multirow}
\usepackage{enumitem}
\usepackage{tikz}
\usetikzlibrary{arrows.meta,positioning,shapes.geometric,fit,calc,matrix}
\usepackage[most]{tcolorbox}
\usepackage{xcolor}
\usepackage{fancyhdr}
\usepackage{hyperref}
\usepackage{microtype}
\usepackage{listings}
\usepackage{pifont}
\setlength{\parindent}{2em}
\setlength{\parskip}{0.32em}
\setlist[itemize]{itemsep=0.18em,topsep=0.28em}
\setlist[enumerate]{itemsep=0.18em,topsep=0.28em}
\hypersetup{colorlinks=true,linkcolor=blue!45!black,urlcolor=blue!50!black}

\definecolor{mainblue}{RGB}{45,86,140}
\definecolor{deepblue}{RGB}{30,62,102}
\definecolor{softblue}{RGB}{238,245,252}
\definecolor{softgreen}{RGB}{238,249,243}
\definecolor{softorange}{RGB}{253,246,233}
\definecolor{softred}{RGB}{253,239,239}
\definecolor{softgray}{RGB}{246,247,249}
\definecolor{deepgreen}{RGB}{35,105,75}
\definecolor{deeporange}{RGB}{165,98,22}
\definecolor{deepred}{RGB}{150,55,55}
\definecolor{purple}{RGB}{94,66,140}
\definecolor{softpurple}{RGB}{245,241,251}
\definecolor{teal}{RGB}{38,112,115}
\definecolor{softteal}{RGB}{238,248,248}
\definecolor{gold}{RGB}{156,120,28}
\definecolor{softgold}{RGB}{252,249,235}

\pagestyle{fancy}
\fancyhf{}
\lhead{408 OS 专题}
\rhead{I/O：请求、搬运、等待与完成}
\cfoot{\thepage}
\renewcommand{\headrulewidth}{0.4pt}
\setlength{\headheight}{14pt}

\newtcolorbox{corebox}[1][]{colback=softblue,colframe=mainblue,title=#1,fonttitle=\bfseries,boxrule=0.8pt,arc=2mm}
\newtcolorbox{exam}[1][]{colback=softorange,colframe=deeporange,title=#1,fonttitle=\bfseries,boxrule=0.8pt,arc=2mm}
\newtcolorbox{warn}[1][]{colback=softred,colframe=deepred,title=#1,fonttitle=\bfseries,boxrule=0.8pt,arc=2mm}
\newtcolorbox{eng}[1][]{colback=softgreen,colframe=deepgreen,title=#1,fonttitle=\bfseries,boxrule=0.8pt,arc=2mm}
\newtcolorbox{method}[1][]{colback=softgray,colframe=black!45,title=#1,fonttitle=\bfseries,boxrule=0.6pt,arc=2mm}
\newtcolorbox{bridge}[1][]{colback=softpurple,colframe=purple,title=#1,fonttitle=\bfseries,boxrule=0.8pt,arc=2mm}
\newtcolorbox{statebox}[1][]{colback=softteal,colframe=teal,title=#1,fonttitle=\bfseries,boxrule=0.8pt,arc=2mm}
\newtcolorbox{history}[1][]{colback=softgold,colframe=gold,title=#1,fonttitle=\bfseries,boxrule=0.8pt,arc=2mm}

\newcommand{\kw}[1]{\textbf{#1}}
\newcommand{\code}[1]{\texttt{#1}}
\newcommand{\arr}{\;\longrightarrow\;}
\newcommand{\yes}{\textcolor{deepgreen}{\ding{51}}}
\newcommand{\no}{\textcolor{deepred}{\ding{55}}}

\lstset{
  basicstyle=\ttfamily\small,
  columns=fullflexible,
  frame=single,
  rulecolor=\color{black!25},
  breaklines=true,
  showstringspaces=false,
  aboveskip=0.6em,
  belowskip=0.6em
}

\title{\vspace{-1.2cm}\textbf{I/O 系统：请求、搬运、等待与完成}\\[0.35em]
\large 从系统调用、轮询、中断、DMA 到缓冲、页缓存、假脱机与磁盘调度}
\author{408 专题手册 · 操作系统 I/O 篇}
\date{}

\begin{document}
\maketitle

\begin{corebox}[核心模型：异步设备怎样接入指令执行过程]
\begin{center}
\Large\bfseries CPU 与进程按指令流推进，而设备速度各异、独立工作、完成时间不确定。\par
I/O 系统的任务，是把“请求、等待、搬运、完成、缓存、排队”组织成一个可控的状态机。
\end{center}

分析主线：
\[
\boxed{
\begin{aligned}
\text{Request}&\arr\text{Submit}\arr\text{Wait / Continue}\arr\text{Transfer}\\
&\arr\text{Complete}\arr\text{Wake / Notify}\arr\text{Return}
\end{aligned}}
\]

做任何 I/O 题，都同时追踪三条线：
\[
\boxed{\text{Control Flow}}
\qquad
\boxed{\text{Data Flow}}
\qquad
\boxed{\text{Process State Flow}}
\]
\end{corebox}

\begin{method}[术语预备：设备怎样和 CPU 交换状态]
\begin{itemize}[leftmargin=1.5em]
  \item \kw{设备 (device)}：执行输入、输出或存储工作的外部硬件；\kw{控制器 (controller)}：把 CPU 可见寄存器/队列和物理设备动作连接起来的硬件接口。
  \item \kw{设备驱动 (driver)}：内核中理解某类控制器寄存器、命令格式和完成通知的软件模块；应用程序通常不直接操作设备寄存器。
  \item \kw{I/O 请求 (request)}：内核为一次读、写或控制操作维护的对象，包含目标、缓冲区、长度、状态和完成后的处理信息。
  \item \kw{轮询 (polling)}：CPU 反复读取状态寄存器或完成队列来检查设备是否完成；\kw{中断驱动}则由设备主动通知 CPU。
  \item \kw{ISR}：Interrupt Service Routine，中断服务程序；它先确认和处理硬件事件，再更新内核状态，必要时唤醒等待者。
  \item \kw{DMA}：Direct Memory Access，设备或 DMA 控制器在 CPU 建立描述符后，直接在设备与内存之间搬运数据，CPU 不必逐字节参与。
  \item \kw{缓冲 (buffer)}：用于吸收生产和消费速率差异、暂存正在搬运的数据；\kw{缓存 (cache)}：保存可由其他副本重新得到的数据以减少慢路径访问。
  \item \kw{假脱机 (Spooling)}：把面向独占设备的请求先写入磁盘队列，再由后台服务顺序提交，使多个程序可以共享打印机等设备。
\end{itemize}
因此“设备完成”“DMA 搬完”“ISR 执行”和“进程被唤醒”是四个不同事件，题目必须按时间顺序分别判断。
\end{method}

\tableofcontents
\newpage

\section{第 0 章：为什么 I/O 必须成为独立子系统}

\subsection{五个根本矛盾}
I/O 不是“CPU 外面挂几个设备”这么简单。它之所以形成完整子系统，是因为至少同时存在五种矛盾。

\begin{center}

\begingroup
\small
\renewcommand{\arraystretch}{1.25}
\setlength{\tabcolsep}{6pt}
\begin{tabularx}{\linewidth}{>{\bfseries}p{0.22\linewidth} X X}
\toprule
根本矛盾 & 如果不解决会怎样 & 自然产生的机制 / 数据结构\\
\midrule
速度不匹配 & CPU 被慢设备拖住；小粒度请求频繁穿越慢路径 & Buffer、Cache、Queue、Read-ahead\\
完成时间不确定 & CPU 不知道设备什么时候完成，只能反复询问 & Interrupt、Completion、Wait/Wakeup\\
搬运成本高 & CPU 花大量指令在设备与内存之间重复搬数据 & DMA (Direct Memory Access)、I/O Channel\\
设备接口异构 & 上层必须理解每种控制器寄存器和协议 & Driver 驱动程序、Device-independent I/O\\
资源有限/独占 & 多个进程争一个打印机、磁盘队列或设备通道 & Scheduling 磁盘调度、Spooling、Device Allocation\\
\bottomrule
\end{tabularx}
\endgroup
\end{center}

\begin{bridge}[I/O 的四次“卸载”]
把 I/O 技术按“CPU/进程被什么拖住”来理解，会比按年代背诵更稳定：
\[
\boxed{
\begin{aligned}
\text{Polling}\to\text{Interrupt} &:\quad \text{卸载等待}\\
\text{PIO}\to\text{DMA} &:\quad \text{卸载搬运}\\
\text{Device-specific}\to\text{Driver} &:\quad \text{卸载设备知识}\\
\text{Direct coupling}\to\text{Buffer/Cache/Spool} &:\quad \text{卸载时间与速率耦合}
\end{aligned}}
\]
\end{bridge}

\subsection{三条并行线：做综合题时必须同时看}

\paragraph{控制线 Control Flow}
谁当前拥有 CPU？控制权怎样进入内核，又怎样从设备完成事件回到内核？
\[
\text{User}\arr\text{Syscall}\arr\text{Kernel}\qquad
\text{Device}\arr\text{Interrupt}\arr\text{Kernel}
\]

\paragraph{数据线 Data Flow}
数据到底从哪里移动到哪里？谁执行搬运？
\[
\text{Device}\xrightarrow{\text{DMA}}\text{RAM/页缓存}
\xrightarrow{\text{copy}}\text{User Buffer}
\]

\paragraph{状态线 Process State Flow}
请求发出以后，原执行流继续、忙等、阻塞还是稍后被唤醒？
\[
\text{Running}\arr\text{Blocked}\arr\text{Ready}\arr\text{Running}
\]

\begin{warn}[三条线最常见的混淆]
“DMA 完成”描述数据搬运；“设备发中断”描述完成通知；“ISR 把等待进程唤醒”描述进程状态变化。三者通常前后相连，但绝不是同一件事。\kw{I/O 完成也不等于原进程立即 Running}：通常只是 Blocked $\to$ Ready，是否马上获得 CPU 仍由调度器决定。
\end{warn}

\section{第 1 章：I/O 请求路径与完成路径}

\subsection{请求怎样进入内核}
用户进程不能直接随意编程设备控制器、修改中断状态或操作关键特权资源。因此 I/O 请求通常通过受控系统调用入口进入内核：
\[
\boxed{\text{User Request}\arr\text{System Call}\arr\text{Kernel I/O Path}}
\]
例如 \code{read(fd, buf, n)} 的语义入口属于用户主动发起的同步控制转移；至于系统调用具体使用 \code{int}、\code{syscall}、\code{svc}、\code{ecall} 等指令，是 ISA/ABI 细节。

\subsection{完成怎样重新进入内核}
设备启动以后通常独立工作。它完成后通过中断、完成队列等机制通知 CPU/内核：
\[
\boxed{\text{Device Completion}\arr\text{Interrupt/Event}\arr\text{Kernel Completion Path}}
\]
因此 I/O 最基本的对偶是：
\[
\boxed{\text{System Call = Request enters kernel}}
\qquad
\boxed{\text{Interrupt = asynchronous event enters kernel}}
\]

\subsection{中断、同步异常与“Trap”术语边界}
\begin{center}

\begingroup
\small
\renewcommand{\arraystretch}{1.25}
\setlength{\tabcolsep}{6pt}
\begin{tabularx}{\linewidth}{>{\bfseries}p{0.18\linewidth} X X X}
\toprule
维度 & 外部中断 Interrupt & 同步异常 Exception & 系统调用/受控陷入 \\
\midrule
触发关系 & 与当前指令语义无关，外部异步到来 & 由当前指令执行引起 & 当前指令有意请求特权服务 \\
典型来源 & 时钟、网卡、磁盘完成、键盘 & Page Fault、除零、非法指令、保护错误 & \code{syscall}/\code{svc}/\code{ecall}/历史 \code{int} \\
核心判断 & 异步事件 & 同步事件 & 同步且主动 \\
处理后动作 & 返回原执行流或调度其他任务 & 修复后重启、向进程报告错误或终止 & 返回系统调用下一阶段或阻塞/调度 \\
\bottomrule
\end{tabularx}
\endgroup
\end{center}

\begin{exam}[408 语言与体系结构严格术语要分层]
教材常把“用户程序主动执行特殊指令进入内核”称为\kw{陷入 Trap}，这在理解系统调用时完全可用；但不要把“Interrupt/Trap/Exception”当成所有 ISA 都完全一致的三分法。判断时直接检查：事件与当前指令同步还是异步、程序是否主动触发、处理后能否继续，以及体系结构规定从哪里重新执行。
\end{exam}

\subsection{Mode Switch 不等于 Context Switch}
一次系统调用或中断首先意味着：
\[
\text{User/Current Context}\arr\text{Kernel Handler}
\]
它\kw{可能仍然属于同一个进程/线程的执行上下文}。只有当内核进一步决定把 CPU 交给另一个可运行实体时，才发生：
\[
\boxed{\text{Context Switch}}
\]
所以：
\[
\boxed{\text{Interrupt/Syscall} \not\Rightarrow \text{Context Switch}}
\]

\section{第 2 章：从最原始的设备交互开始}

\subsection{控制器是 CPU 与物理设备之间的硬件接口}
典型设备控制器至少暴露三类逻辑信息：
\begin{itemize}
  \item \kw{Status}：忙/空闲、错误、完成状态；
  \item \kw{Command/Control}：读、写、启动、复位等命令；
  \item \kw{Data/Descriptor}：数据寄存器、DMA 描述符地址、队列门铃等。
\end{itemize}
CPU/驱动通过端口 I/O 或 MMIO 等方式访问这些接口。

\begin{eng}[工程边界：MMIO 比“某个固定端口”更通用]
现代设备大量通过 Memory-Mapped I/O 暴露寄存器，CPU 发出的特定地址访问会被互连/总线送到设备而非普通 DRAM。408 题目更关注“控制/状态/数据寄存器 + 驱动访问”的逻辑模型，不必把某一总线实现背成唯一形式。
\end{eng}

\subsection{程序查询方式：最简单，也最浪费 CPU}
最直接方案：
\begin{lstlisting}
start_device();
while (status() == BUSY) {
    ; // busy waiting
}
read_result();
\end{lstlisting}

CPU 在循环中反复读取设备状态寄存器：
\[
\boxed{\text{CPU repeatedly pays to discover that nothing changed}}
\]
如果设备操作远慢于 CPU，绝大多数检查都是无效工作。

\begin{exam}[程序查询 I/O 的核心特征]
\begin{itemize}
  \item CPU 主动反复查询设备状态；
  \item CPU 与设备工作可以在时间上交错，但等待期间 CPU 被忙等占住；
  \item 数据搬运通常仍由 CPU 指令参与；
  \item 实现简单，适合极短等待、简单控制器或某些低延迟特殊场景。
\end{itemize}

\begin{method}[设备、控制器、接口与端口四层辨析]
物理设备执行动作；设备控制器解释命令并管理数据/状态；I/O 接口规定 CPU/总线怎样访问控制器；I/O 端口是命令、状态、数据寄存器的可寻址入口。端口可采用独立编址（专用 I/O 指令）或统一编址/内存映射 I/O（普通 load/store 地址空间），两者改变的是访问指令与地址空间，不改变设备寄存器必须受内核权限保护的事实。
\end{method}
\end{exam}

\section{第 3 章：中断驱动 I/O——把“等待”从 CPU 上卸载}

\subsection{为什么中断自然产生}
如果设备完成时间不可预测，最合理的协议不是 CPU 一直问，而是：
\[
\boxed{\text{Start}\arr\text{Do Other Work}\arr\text{Device Notifies Completion}}
\]
于是：
\[
\boxed{\text{Polling}\to\text{Interrupt}=\text{Offload Waiting}}
\]

\subsection{一次中断驱动 I/O 的完整状态路径}
假设进程 A 发起阻塞式设备读取：
\[
A:\ \text{Running}
\arr \text{syscall}
\arr \text{driver starts I/O}
\arr \text{Blocked}
\]
CPU 转去运行 B。设备完成以后：
\[
\text{Device}\arr\text{IRQ}\arr\text{ISR/Completion}\arr A:\text{Ready}
\]
未来调度器选中 A：
\[
A:\text{Ready}\arr\text{Running}\arr\text{return to user}
\]

\subsection{中断本身也有成本}
中断不是“免费通知”。至少可能包含：
\begin{itemize}
  \item 保存必要体系结构状态并进入内核；
  \item 执行 ISR/完成处理；
  \item cache/TLB/流水线局部扰动；
  \item 可能产生调度和上下文切换。
\end{itemize}
因此，若事件极短、马上完成，短时间 polling/spin 可能比睡眠再唤醒更经济。

\begin{bridge}[与并发专题同构：Polling vs Interrupt = Spin vs Sleep]
两组问题都要区分“谁等待”和“谁搬运数据”：
\[
\boxed{\text{短等待：支付 CPU 时间，避免切换成本}}
\]
\[
\boxed{\text{长等待：放弃 CPU，支付阻塞/唤醒/调度成本}}
\]
所以不存在脱离等待时间与负载的“interrupt 永远优于 polling”。
\end{bridge}

\begin{eng}[现代工程：Interrupt moderation / batching]
高速网卡、NVMe 等设备可能每秒完成海量小请求。若每个完成都单独中断，CPU 会被中断风暴淹没。因此现实系统常批量完成、合并中断或使用 polling 与 interrupt 的混合模式。它们改变的是工程 快路径，不改变 408 的核心模型：\kw{完成事件必须被系统发现，并映射成软件状态更新}。
\end{eng}

\section{第 4 章：DMA 与 I/O 通道——把“搬运”进一步卸载}

\subsection{Interrupt 只解决“等”，没有自动解决“搬”}
若每个数据字节仍按：
\[
\text{Device}\arr\text{CPU Register}\arr\text{Memory}
\]
移动，那么 CPU 虽然不用忙等，却还要执行大量搬运指令。因此进一步引入：
\[
\boxed{\text{PIO}\to\text{DMA}=\text{Offload Data Movement}}
\]

\subsection{DMA 四阶段模型}
\begin{center}

\begingroup
\small
\begin{tikzpicture}[>=Stealth, node distance=1.8cm and 3.2cm]
\tikzset{
  n1/.style={draw=mainblue, ultra thick, fill=softblue, rounded corners=2.5mm, align=center, minimum width=2.8cm, minimum height=0.95cm, font=\small\bfseries},
  n2/.style={draw=purple, ultra thick, fill=softpurple, rounded corners=2.5mm, align=center, minimum width=2.8cm, minimum height=0.95cm, font=\small\bfseries},
  n3/.style={draw=deeporange, ultra thick, fill=softorange, rounded corners=2.5mm, align=center, minimum width=2.8cm, minimum height=0.95cm, font=\small\bfseries},
  n4/.style={draw=deepgreen, ultra thick, fill=softgreen, rounded corners=2.5mm, align=center, minimum width=2.8cm, minimum height=0.95cm, font=\small\bfseries}
}
\node[n1] (setup) {1. Setup 驱动配置\\[2pt]\small 缓冲区与 DMA 描述符};
\node[n2, right=of setup] (transfer) {2. Transfer DMA 传输\\[2pt]\small 设备与 RAM 直接读写};
\node[n3, right=of transfer] (irq) {3. IRQ 触发中断\\[2pt]\small 硬件告知传输完成};
\node[n4, below=1.8cm of irq] (cleanup) {4. Cleanup 请求收尾\\[2pt]\small 唤醒进程 / 回收资源};

\draw[->, ultra thick, deepblue] (setup) -- node[above=3pt, font=\small\bfseries] {启动硬件 DMA} (transfer);
\draw[->, ultra thick, purple] (transfer) -- node[above=3pt, font=\small\bfseries] {传输结束} (irq);
\draw[->, ultra thick, deeporange] (irq) -- node[right=4pt, font=\small\bfseries] {进入 ISR} (cleanup);
\draw[->, ultra thick, deepgreen] (cleanup.south) -- ++(0,-1.0cm) -| node[pos=0.25, below=3pt, font=\small\bfseries] {解阻塞返回} (setup.south);
\end{tikzpicture}
\endgroup
\end{center}

经典 DMA 控制器至少要表达四类状态：命令/状态寄存器 CR、内存地址寄存器 MAR、数据暂存寄存器 DR 和计数器 DC。CPU 预处理时给出方向、起始地址与长度；DMA 传送时按单一 MAR 前进、DC 递减；完成后以一次中断报告整块结果。因而经典单组 MAR/DC 只能描述一段连续内存，scatter-gather 或通道程序属于更强的描述能力。所谓“数据不经过 CPU”只指不经 CPU 通用寄存器，DMA 仍会争用内存总线并要求映射在传输期间有效。

\begin{statebox}[DMA 的正确一句话]
DMA \kw{不是“完全不需要 CPU”}，而是把大块数据搬运的逐字节/逐字指令参与从 CPU 上卸载。CPU 仍负责建立请求、编程/提交描述符、处理中断或完成队列，并管理相关内存与同步状态。
\end{statebox}

\subsection{DMA 仍然消耗系统资源}
即使 CPU 不逐字搬数据，传输仍然会占用：
\begin{itemize}
  \item 内存带宽；
  \item 系统互连/总线带宽；
  \item 设备队列与控制器资源；
  \item 某些平台上的 cache coherence / IOMMU 映射资源。
\end{itemize}
所以：
\[
\boxed{\text{DMA}\neq\text{zero cost}}
\]

\begin{eng}[工程边界：DMA address 不一定等于 CPU physical address]
现代平台常存在 IOMMU 或总线地址翻译。驱动通常通过 DMA API 得到设备可使用的 DMA address；设备使用该地址访问内存，必要时再由 IOMMU 翻译到真实系统内存。不要把“DMA 直接写物理地址”当成所有平台的绝对实现。
\end{eng}

\subsection{DMA 映射生命周期：设备看见的内存必须在传输期间保持稳定}

来源笔记中“DMA 期间不能让页框被回收或改作他用”的直觉是正确的，但不能把它简化为“所有体系结构的 PTE 都有一个 pinned bit”。稳定的接口契约是：
\[
\boxed{
\begin{aligned}
\text{prepare buffer}
&\arr\text{establish DMA-visible mapping}
\arr\text{transfer}\\
&\arr\text{complete/synchronize}
\arr\text{unmap or release}
\end{aligned}
}
\]

在这段生命周期内必须保持：
\begin{itemize}
  \item 设备使用的 DMA address 仍映射到预期内存，相关页面/缓冲不能被回收并复用于别的对象；
  \item 传输方向与 ownership 清楚，CPU 和设备何时可以读取或修改缓冲必须符合 DMA API 与一致性要求；
  \item mapping/pinning 与 unmapping/unpinning 成对出现，失败路径也要回收资源。
\end{itemize}

若用户缓冲在虚拟地址上连续、物理上离散，可以使用 scatter--gather 描述多个片段；设备能力或地址限制不允许时，可能需要 bounce buffer。二者解决的是“设备怎样看到可访问的地址序列”，并不取消映射生命周期。

\begin{warn}[三个不能绝对化的旧稿说法]
DMA 设备不一定只看到裸 CPU physical address（可能经 IOMMU）；pinning 的实现状态不一定存放在 PTE 固定位；内核分配的普通内存也不能仅凭分配函数名称推断 DMA 可达、物理连续或无需映射。考试题按给定教材模型作答，工程边界以平台 DMA API 为准。
\end{warn}

\subsection{408 经典：DMA 的几种总线使用方式}
不同教材术语略有差异，但核心权衡是“DMA 连续占用互连多长时间”。
\begin{itemize}
  \item \kw{Cycle Stealing（周期挪用）}：DMA 每次取得一个/少量总线周期，与 CPU 交替使用；
  \item \kw{Burst/Block Transfer（突发/块传送）}：DMA 一次连续传送较大数据块，吞吐高但 CPU 可能更久拿不到总线；
  \item \kw{Transparent DMA（透明方式）}：尽量利用 CPU 不访问存储器/总线的空隙传输，干扰小但实现和可用时机受限。
\end{itemize}
做题时不要只背名称，要问：\kw{谁占用总线？持续多久？CPU 被延迟多少？总体吞吐如何？}

\subsection{I/O Channel：比 DMA 更进一步的控制卸载}
在经典大型机/教材模型中，I/O 通道可以执行专门的通道程序，管理一串 I/O 操作：
\[
\boxed{\text{CPU issues high-level I/O job}\arr\text{Channel executes I/O program}}
\]
相较 DMA 主要卸载数据搬运，通道还卸载更多\kw{设备控制与 I/O 操作序列组织}。408 若出现“通道”概念，应抓住这个层级区别，而不是把它简单等同于 DMA 控制器。

\subsection{四种经典 I/O 控制方式总对比}
\begingroup
\small
\renewcommand{\arraystretch}{1.25}
\setlength{\tabcolsep}{6pt}
\begin{tabularx}{\linewidth}{>{\bfseries}p{0.18\linewidth} p{0.2\linewidth} p{0.22\linewidth} X X}
\toprule
方式 & CPU 是否忙等 & 数据主要由谁搬 & 完成/控制粒度 & 核心评价 \\
\midrule
程序查询 & 是 & CPU / PIO & CPU 反复检查，常按小单位推进 & 硬件简单，CPU 浪费最大 \\
中断驱动 & 否（可去做别的） & 经典模型中 CPU 仍较多参与 & 设备准备/完成时中断 CPU & 卸载等待，但中断和搬运仍有开销 \\
DMA & 否 & DMA 引擎/设备 $\leftrightarrow$ 内存 & CPU 配置块传输，完成后再通知 & 大幅减少 CPU 搬运参与 \\
I/O 通道 & 否 & 通道执行 I/O 程序并组织设备操作 & CPU 提交更高层 I/O 任务 & 进一步卸载控制与序列组织 \\
\bottomrule
\end{tabularx}
\endgroup

\begin{exam}[最稳的比较方法：不要背“谁最好”]
对每种控制方式统一问四件事：\kw{CPU 是否忙等？CPU 是否逐单位搬数据？一次需要多少次 CPU/中断参与？额外硬件复杂度有多大？} 技术演进的主方向是降低 CPU 对慢 I/O 的细粒度参与，但硬件和协议复杂度随之上升。
\end{exam}

\subsection{经典 I/O 通道类型：按“设备共享通道的方式”理解}
部分 408/教材体系还会区分：
\begin{itemize}
  \item \kw{字节多路通道}：多个低速设备按字节/小单位交叉共享通道，强调多路并发；
  \item \kw{数组选择通道}：一次选择一个高速设备，连续完成一批数据传输后再换设备，吞吐高但通道独占时间较长；
  \item \kw{数组多路通道}：兼顾块传输与多设备交叉，在多个高速设备的块操作之间复用通道。
\end{itemize}
不要孤立背名称。比较三类通道时，依次检查：通道一次交给谁、占用多久、每次传多少数据、能否并发服务多个设备。

\section{第 5 章：Driver 与 Device-independent I/O——把设备知识隔离}

\subsection{为什么必须有驱动}
若文件系统、网络栈、应用层都直接理解每个设备寄存器：
\[
\text{Upper Layer}\times\text{Every Device Model}
\]
系统复杂度会爆炸。驱动建立：
\[
\boxed{\text{Common OS Request}\arr\text{Device-specific Commands}}
\]
于是：
\[
\boxed{\text{Device-specific Knowledge}\to\text{Driver}}
\]

\subsection{经典 I/O 软件层次}
\begin{center}

\begingroup
\small
\begin{tikzpicture}[>=Stealth, node distance=0.8cm]
\tikzset{
  l1/.style={draw=mainblue, ultra thick, fill=softblue, rounded corners=2.5mm, align=center, minimum width=14.0cm, minimum height=0.95cm, font=\small\bfseries},
  l2/.style={draw=purple, ultra thick, fill=softpurple, rounded corners=2.5mm, align=center, minimum width=14.0cm, minimum height=0.95cm, font=\small\bfseries},
  l3/.style={draw=deeporange, ultra thick, fill=softorange, rounded corners=2.5mm, align=center, minimum width=14.0cm, minimum height=0.95cm, font=\small\bfseries},
  l4/.style={draw=deepgreen, ultra thick, fill=softgreen, rounded corners=2.5mm, align=center, minimum width=14.0cm, minimum height=0.95cm, font=\small\bfseries},
  l5/.style={draw=deepred, ultra thick, fill=softred, rounded corners=2.5mm, align=center, minimum width=14.0cm, minimum height=0.95cm, font=\small\bfseries}
}
\node[l1] (u) {1. 用户级 I/O 软件 / 库函数 / 应用 (User-level I/O)};
\node[l2, below=of u] (ind) {2. 设备独立性 I/O 层：统一命名、保护、缓冲、错误处理与设备分配};
\node[l3, below=of ind] (drv) {3. 设备驱动程序：翻译特定命令、编程控制寄存器与 DMA 描述符};
\node[l4, below=of drv] (int) {4. 中断/完成处理程序：响应硬件事件、解阻塞与唤醒等待线程};
\node[l5, below=of int] (hw) {5. 硬件设备与控制器：执行物理读写与数据传输};

\draw[->, ultra thick, deepblue] (u) -- node[right=4pt, font=\small\bfseries] {系统调用请求} (ind);
\draw[->, ultra thick, purple] (ind) -- node[right=4pt, font=\small\bfseries] {设备无关标准接口} (drv);
\draw[->, ultra thick, deeporange] (drv) -- node[right=4pt, font=\small\bfseries] {驱动特定指令与描述符} (int);
\draw[<->, ultra thick, deepgreen] (int) -- node[right=4pt, font=\small\bfseries] {总线读写 / 中断通知} (hw);
\end{tikzpicture}
\endgroup
\end{center}

\begin{exam}[为什么“设备独立层”很重要]
它回答的不是“某个磁盘怎么转”，而是“无论设备型号如何，上层都希望拥有统一的 open/read/write、保护、命名、缓冲、错误返回和资源分配语义”。这与 VFS、网络 socket 抽象、虚拟内存中的地址间接层是同一种复杂度控制思想。
\end{exam}

\subsection{设备标识也要分层：逻辑名字、设备号与驱动实例}

设备独立性不表示内核不再区分具体设备，而是把“上层怎样命名”与“哪个驱动处理哪个实例”隔开。在类 UNIX 教学模型中，设备结点可携带
\[
\boxed{device\ id=(major,minor)}.
\]
其中 major 通常选择驱动/设备类别的处理入口，minor 再由该驱动解释为具体实例、分区或子设备。稳定链条是：
\[
\text{logical/device name}
\arr\text{device identity}
\arr\text{driver dispatch}
\arr\text{concrete instance}.
\]

\begin{warn}[不要把两组“主从”混在一起]
major/minor number 是 OS 的分类与实例标识；总线语境中的 bus master/slave 描述谁发起、谁响应一次传输。名称相似，但对象、关系和可推出结论都不同。
\end{warn}

\subsection{Cache vs Buffer：\texorpdfstring{$A\neq B$}{A != B} 核心功能辨析}

考场与实际系统中，缓存 (Cache) 与缓冲 (Buffer) 经常被混淆。可按数据取走后是否保留来区分：
\[
\boxed{\text{Cache: “下次可能还要，所以留下副本”}}
\]
\[
\boxed{\text{Buffer: “现在速度/粒度对不上，所以暂存”}}
\]

\begingroup
\small
\renewcommand{\arraystretch}{1.25}
\setlength{\tabcolsep}{6pt}
\begin{tabularx}{\linewidth}{>{\bfseries}p{0.18\linewidth} X X}
\toprule
对比维度 & 缓存 (Cache) & 缓冲 (Buffer)\\
\midrule
核心目的 & 利用时间/空间局部性，**减少未来的重复访问延迟** & 协调两端**速率 mismatch、粒度差异与同步交错**\\
数据独立性 & 强调是物理主存/后备对象的**副本 (Copy)** & 暂存**传输中**的数据，不一定强调副本概念\\
典型场景 & Page Cache、TLB、指令/数据 Cache & 磁盘写缓冲、网络 Socket 发送/接收 Buffer、打印缓冲\\
兼具兼顾场景 & \multicolumn{2}{p{0.75\linewidth}}{真实 OS 内存区（如 Linux Page Cache）可**同时承担双重职责**：既作为 I/O 磁盘读写的速度/粒度缓冲，又保留映射作为后续读访问的缓存副本。}\\
\bottomrule
\end{tabularx}
\endgroup

\subsection{字符设备与块设备：不同工作负载产生不同接口}
\begingroup
\small
\renewcommand{\arraystretch}{1.25}
\setlength{\tabcolsep}{6pt}
\begin{tabularx}{\linewidth}{>{\bfseries}p{0.2\linewidth} X X}
\toprule
维度 / 特性 & 字符/流式设备（经典抽象） & 块设备（经典抽象） \\
\midrule
数据组织形式 & 连续字节/字符流，强顺序消费与单向流动 & 固定大小块（扇区/扇区组），支持随机寻址 \\
典型代表设备 & 键盘、串口、鼠标、部分传感器 & 机械硬盘 HDD、固态硬盘 SSD、块存储 \\
OS 系统重点 & 流控、字符缓冲、终端行规程、事件到达 & 块队列、磁盘调度、页缓存 缓存、块映射 \\
\bottomrule
\end{tabularx}
\endgroup
实际设备/驱动接口比这更丰富，但这一区分有助于理解为什么磁盘能形成 页缓存 与块层，而键盘不适合用“随机访问块”理解。

\subsection{网络设备：流量包与设备队列的接口层}
网络接口既不是普通字符流，也不是可任意寻址的磁盘块。驱动通常面对接收/发送描述符环、数据包缓冲区和设备完成事件；上层通过统一的网络设备 API 把包提交给协议栈，协议栈再把端点、顺序、重传等语义交给网络协议 Owner。本册只追踪 I/O 侧的对象：
\[
\text{RX/TX descriptor}\to\text{buffer ownership}\to\text{device progress}\to\text{completion}\to\text{wake/notify}.
\]
发送时必须明确 CPU/驱动何时放弃缓冲所有权，接收时必须明确设备何时把填好的缓冲交还内核；不能把“收到一个中断”直接当成应用已经收到完整消息。网络设备可用 DMA 批量搬运，也可在非阻塞 API 下立即返回“当前无包/队列满”，这与字符设备的流式读取共享控制路径，但数据单元和排队策略不同。

\subsection{I/O 保护：为什么设备寄存器和 DMA 也需要权限边界}
I/O 设备能够读写外部世界，某些设备还能直接 DMA 到内存，因此它们天然具有强权限。系统必须控制：
\begin{itemize}
  \item 用户态不能任意执行特权 I/O 指令或修改关键控制器状态；
  \item MMIO 区域不能任意映射给不可信进程；
  \item DMA 缓冲和 DMA 地址必须由内核/驱动建立合法映射；
  \item IOMMU（若存在）可以限制设备可 DMA 的地址范围，降低设备或驱动错误破坏任意内存的风险。
\end{itemize}
因此 I/O 保护仍然符合本系列总原则：
\[
\boxed{\text{Who can command which device to access which memory?}}
\]

\subsection{错误与完成码：I/O 慢路径 不只有“成功”}
设备完成事件可能携带：成功、短传输、超时、校验/介质错误、设备离线等状态。内核需要决定：
\[
\text{Retry?}\quad\text{Fail?}\quad\text{Remap/Recover?}\quad\text{Wake which waiter?}
\]
因此一个稳健的 I/O 状态机不仅有 Request $\to$ Complete，还要允许：
\[
\boxed{\text{Request}\arr\text{Pending}\arr\{\text{Success},\text{Retry},\text{Error},\text{Cancel}\}}
\]
408 通常不深入驱动错误恢复细节，但做综合题时必须意识到“中断到来”只说明有事件，\kw{不自动等于本次 I/O 成功完成}。

\section{第 6 章：四个维度必须拆开——Blocking、Async、Interrupt、DMA}

\subsection{为什么这几个词特别容易乱}
因为它们回答的是\kw{不同问题}：

\begingroup
\small
\renewcommand{\arraystretch}{1.25}
\setlength{\tabcolsep}{6pt}
\begin{tabularx}{\linewidth}{>{\bfseries}p{0.25\linewidth} X X}
\toprule
维度概念 & 需要回答的问题 & 两端概念与核心对比 \\
\midrule
调用线程是否等待 & 调用暂时不能完成时，进程/线程现在怎么办？ & Blocking (阻塞) / Nonblocking (非阻塞) \\
完成语义如何交付 & 调用返回时结果是否已经完成？完成以后怎样通知？ & Synchronous (同步) / Asynchronous (异步) \\
设备完成如何被发现 & CPU 主动反复查询，还是设备主动中断通知？ & Polling (程序查询) / Interrupt (中断驱动) \\
数据主要由谁搬运 & CPU 指令逐步搬运，还是设备/DMA 引擎搬运？ & PIO (Programmed I/O) / DMA \\
\bottomrule
\end{tabularx}
\endgroup

因此完全可能有：
\[
\boxed{\text{Blocking + Interrupt-driven + DMA}}
\]
普通阻塞磁盘读取就是典型直觉模型：线程阻塞；设备 DMA；完成以后中断；最后线程被唤醒。

\subsection{阻塞与非阻塞只看“当前线程”}
\paragraph{Blocking}
条件暂不满足时：
\[
\text{Running}\arr\text{Blocked}
\]
CPU 可运行别的线程。

\paragraph{Nonblocking}
条件暂不满足时立即返回：
\[
\text{return EAGAIN/partial/none}
\]
应用稍后重试或结合事件机制。

\subsection{同步与异步要看“完成语义”}
不要把“用了中断”直接等同于“异步 I/O”。设备内部即使由中断完成，用户线程仍可能执行一个\kw{同步阻塞}系统调用；异步接口则允许请求先返回，完成以后再通过 completion queue、callback、signal/event 等方式交付结果。

\begin{method}[四维判断法]
遇到任何 I/O API/题目，不问“它到底是同步还是中断”，而依次问：
\begin{enumerate}
  \item 调用线程等待吗？
  \item 系统调用返回时数据完成了吗？
  \item 设备完成事件是谁发现的？
  \item 大块数据是谁搬的？
\end{enumerate}
四个问题独立作答，概念自然不会打结。
\end{method}

\section{第 7 章：Buffer、Cache、Spooling——看起来都“存一下”，目的完全不同}

\begingroup
\small
\renewcommand{\arraystretch}{1.25}
\setlength{\tabcolsep}{6pt}
\begin{tabularx}{\linewidth}{>{\bfseries}p{0.18\linewidth} p{0.22\linewidth} X X}
\toprule
机制 & 解决的核心根本问题 & 核心动作与策略 & 典型应用场景 \\
\midrule
Buffer & 生产/消费速度或粒度不匹配 & 暂存、吸收突发、平滑速率 & 键盘缓冲区、网络收发缓冲、双缓冲 \\
Cache & 慢对象可能再次被访问 & 保存副本以复用，命中时绕过慢层 & 页缓存、CPU Cache、TLB \\
Spooling & 独占/慢设备与多个任务直接耦合 & 中转存储 + 排队 + 后台服务 & 打印任务、经典输入/输出井 \\
DMA & CPU 逐字搬运过于昂贵 & 设备/控制器直接与内存传输 & 磁盘、网卡、GPU 数据传输 \\
Interrupt & CPU 主动等待过于昂贵 & 事件发生时主动中断通知 CPU & I/O 完成、时钟、外部硬件事件 \\
\bottomrule
\end{tabularx}
\endgroup

\subsection{Buffer：速率匹配与粒度匹配}
Buffer 的主要任务不是“下次再用”，而是让生产者和消费者不必严格同速：
\[
\text{Producer}\arr\boxed{\text{Buffer}}\arr\text{Consumer}
\]
它可以：
\begin{itemize}
  \item 吸收短时突发；
  \item 把很多小请求聚合成大请求；
  \item 允许设备与 CPU 在一段时间内并行；
  \item 减少双方必须同时在线的程度。
\end{itemize}

\subsection{单缓冲、双缓冲与循环缓冲}
\paragraph{单缓冲}
设备填充和进程取用之间仍有阶段性互斥，同一 buffer 不能一边被设备覆盖、一边被进程消费。

\paragraph{双缓冲}
\[
B_0\ \text{被消费时}\quad B_1\ \text{可被设备填充}
\]
如果生产和消费时间接近，可明显提高重叠程度。

\paragraph{循环缓冲}
多个 buffer 形成 ring：
\[
B_0\to B_1\to\cdots\to B_{N-1}\to B_0
\]
这已经自然连接到生产者-消费者问题：数据结构只是容器，正确性仍需要处理“空/满、读写索引、并发更新”。

\begin{center}
\providecolor{themebg}{HTML}{FAFAF7}
\providecolor{themefg}{HTML}{111111}
\providecolor{themegray}{HTML}{666666}
\providecolor{themecurve}{HTML}{1D63B8}
\providecolor{themealert}{HTML}{C53030}
\providecolor{themeamber}{HTML}{B86E00}
\providecolor{themegreen}{HTML}{15803D}
\begin{tikzpicture}[font=\scriptsize,color=themefg,text=themefg,>=Stealth]
\tikzset{
  tbox/.style={draw=themecurve,fill=themecurve!12!themebg,font=\scriptsize\bfseries,text=themecurve,align=center},
  mbox/.style={draw=themeamber,fill=themeamber!12!themebg,font=\scriptsize\bfseries,text=themeamber,align=center},
  cbox/.style={draw=themegreen,fill=themegreen!12!themebg,font=\scriptsize\bfseries,text=themegreen,align=center},
  lbl/.style={font=\tiny\bfseries,fill=themebg,inner sep=1.5pt,rounded corners=1pt},
  note/.style={rounded corners=4pt,draw=themegray!70,fill=themefg!3!themebg,inner xsep=8pt,inner ysep=7pt,align=left}
}

% ── 标题与例题说明 ──
\node[anchor=west,font=\bfseries\Large] at (0,8.8) {单缓冲与双缓冲：$T$（输入）、$M$（传送到用户区）、$C$（CPU处理）时序与稳态周期};
\node[anchor=west,text=themegray,text width=17.5cm,align=left] at (0,8.2)
  {符号约定：$T$ 为外设向系统缓冲区输入一块数据耗时；$M$ 为系统缓冲区向用户区传送耗时；$C$ 为 CPU 计算处理一块数据耗时。};

% ══════════════════════════════════════════════════════════════
% ── Part A: 单缓冲区 (Single Buffer) (y = 4.8 .. 7.4) ──
% ══════════════════════════════════════════════════════════════
\node[font=\bfseries\normalsize,text=themecurve,anchor=west] at (0,7.5) {A. 单缓冲区：系统缓冲区 $M$ 与 $T$ 互斥（串行瓶颈）};

\node[anchor=east,font=\scriptsize\bfseries] at (2.4,6.6) {外设输入/系统缓冲};
\draw[tbox] (2.6,6.25) rectangle (5.6,6.95) node[midway] {$T_1$};
\draw[mbox] (5.6,6.25) rectangle (7.0,6.95) node[midway] {$M_1$};
\draw[tbox] (7.0,6.25) rectangle (10.0,6.95) node[midway] {$T_2$};
\draw[mbox] (10.0,6.25) rectangle (11.4,6.95) node[midway] {$M_2$};

\node[anchor=east,font=\scriptsize\bfseries] at (2.4,5.4) {用户工作区 CPU};
\draw[cbox] (7.0,5.05) rectangle (9.6,5.75) node[midway] {$C_1$};
\draw[cbox] (11.4,5.05) rectangle (14.0,5.75) node[midway] {$C_2$};

% 周期标注
\draw[<->,line width=1.1pt,themealert] (5.6,4.5) -- (10.0,4.5);
\node[font=\scriptsize\bfseries,text=themealert,anchor=north] at (7.8,4.4) {单缓冲稳态周期：$t_{\text{single}} = \max(T, C) + M$};

% ══════════════════════════════════════════════════════════════
% ── Part B: 双缓冲区 (Double Buffer) (y = 1.0 .. 3.6) ──
% ══════════════════════════════════════════════════════════════
\node[font=\bfseries\normalsize,text=themeamber,anchor=west] at (0,3.3) {B. 双缓冲区：缓冲区 1 与缓冲区 2 轮转流水（$T$ 与 $M+C$ 完全重叠）};

\node[anchor=east,font=\scriptsize\bfseries] at (2.4,2.5) {外设连续输入};
\draw[tbox] (2.6,2.15) rectangle (5.6,2.85) node[midway] {$T_1 \to \text{Buf1}$};
\draw[tbox] (5.6,2.15) rectangle (8.6,2.85) node[midway] {$T_2 \to \text{Buf2}$};
\draw[tbox] (8.6,2.15) rectangle (11.6,2.85) node[midway] {$T_3 \to \text{Buf1}$};

\node[anchor=east,font=\scriptsize\bfseries] at (2.4,1.3) {传送与计算 ($M+C$)};
\draw[mbox] (5.6,0.95) rectangle (6.8,1.65) node[midway] {$M_1$};
\draw[cbox] (6.8,0.95) rectangle (8.6,1.65) node[midway] {$C_1$};
\draw[mbox] (8.6,0.95) rectangle (9.8,1.65) node[midway] {$M_2$};
\draw[cbox] (9.8,0.95) rectangle (11.6,1.65) node[midway] {$C_2$};

% 周期标注
\draw[<->,line width=1.1pt,themealert] (5.6,0.4) -- (8.6,0.4);
\node[font=\scriptsize\bfseries,text=themealert,anchor=north] at (7.1,0.3) {双缓冲稳态周期：$t_{\text{double}} = \max(T, M + C)$};

% ── 底部总结与误读警示 ──
\node[note,anchor=west,text width=17.5cm] at (0,-0.8)
  {\textbf{$n$ 块数据处理总时间计算公式：}\\
   $\bullet$ \textbf{单缓冲总耗时}：$T_{\text{total}} = T + (n-1) \times (\max(T, C) + M) + M + C = n \times (\max(T, C) + M) + \min(T, C)$；\\
   $\bullet$ \textbf{双缓冲总耗时}：$T_{\text{total}} = T + M + n \times C$（若 $T < M+C$）或 $n \times T + M + C$（若 $T > M+C$），即 $T_{\text{total}} = (n-1) \times \max(T, M+C) + T + M + C$。};

\node[anchor=west,font=\scriptsize\bfseries,text=themealert,text width=17.5cm,align=left] at (0,-1.9)
  {禁止误读：不要直接死背 $n \times \text{周期}$！必须先画出第 1 块的冷启动前导流水，加上中间 $(n-1)$ 个稳态周期，再加上尾部排空。};
\end{tikzpicture}

\end{center}

\begin{method}[缓冲时间模型：公式必须绑定假设]
设设备把一块送入内核缓冲耗时 $T$，从内核缓冲复制到用户工作区耗时 $M$，CPU 处理该块耗时 $C$。在教材常用流水线假设下：
\[
\boxed{t_{single}=\max(T,C)+M},\qquad
\boxed{t_{double}=\max(T,C+M)}.
\]
单缓冲中 $M$ 期间唯一缓冲仍被占用，设备不能填下一块；双缓冲中设备可填另一块，所以 $M$ 并入 CPU 侧。$n\ge2$ 时稳态吞吐仍由 $\max(T,C+M)$ 决定，多余缓冲主要吸收突发而不提高慢级本身的速度。双缓冲也不必然更快：当 $C+M\ge T$ 时，两种配置可能得到相同周期。若题目允许关系不同，应回到时间线而不是硬套公式。
\end{method}

循环缓冲的两个追赶事件也可直接判瓶颈：输入指针追上消费指针表示空缓冲耗尽，输入方阻塞、计算侧较慢；消费指针追上输入指针表示满缓冲耗尽，消费方阻塞、I/O 侧较慢。缓冲池的资源量信号量必须先于队列互斥锁申请，否则可能抱锁睡眠；这与生产者--消费者中先申请 empty/full 再进入 mutex 的顺序同源。

\subsection{经典缓冲池：不要背三条队列，要追踪 buffer 状态}

教材常把一组公共缓冲区组织成空闲队列、输入队列与输出队列。分析时不要只背队列名称，要跟踪缓冲区在以下状态间怎样移动：
\[
\boxed{
\text{Free}
\xrightarrow{\text{reserve for input}}
\text{Filling}
\xrightarrow{\text{input complete}}
\text{Ready-for-consumer}
\xrightarrow{\text{release}}
\text{Free}
}
\]
输出方向同理，只是 producer/consumer 的角色交换。每次队列迁移必须同时回答：当前 buffer 由谁拥有、内容是否有效、谁可能等待、什么事件触发下一次迁移。

\begin{exam}[缓冲池题的固定检查]
先画 buffer 状态与队列，再模拟“申请、装入、提取、释放”。若没有空缓冲，申请者等待的是 buffer 资源；若目标输入/输出队列为空，消费者等待的是数据条件。两种等待不能只因为都出现阻塞而合并。
\end{exam}

\begin{exam}[缓冲题不要死背总时间公式]
不同教材题目对“设备填充、内核复制、用户处理哪些阶段可以重叠”的假设不同。稳妥做法是：\kw{先画时间线/甘特图，再根据互斥和可重叠关系计算}。公式只在假设一致时成立。
\end{exam}

\section{第 8 章：页缓存——文件 I/O 的内存 快路径}

\subsection{核心对象：缓存的是“文件内容”}
页缓存 可以抽象为：
\[
\boxed{(File,Offset/PageIndex)\arr Cached\ Memory}
\]
它由内核管理，使普通文件 I/O 的常见访问尽量停留在内存，而不是每次都访问底层存储。

\subsection{读取路径：Hit、Miss 与 Read-ahead}
\paragraph{Hit}
\[
\text{read}\arr\text{页缓存 Hit}\arr\text{copy to user}\arr\text{return}
\]
无底层存储读取。

\paragraph{Miss}
\[
\text{页缓存 Miss}\arr\text{allocate/cache page}\arr\text{submit storage I/O}
\]
若调用是阻塞式读取，当前线程可能 Blocked。设备 DMA 把数据放入内存，完成后中断/完成处理唤醒线程，随后复制到用户 buffer。

\paragraph{Read-ahead}
系统观察到顺序/可预测访问时，可能在应用显式请求之前预取后续文件页：
\[
\boxed{\text{Observed Locality}\arr\text{Predict Future File Pages}\arr\text{Asynchronous Prefetch}}
\]
不要记固定“预读 8 页/32KB”等数字；窗口和策略属于实现。

\subsection{写入路径：Buffered Write 与 Dirty Page}
普通 buffered write 可抽象成：
\[
\text{User Buffer}\arr\text{页缓存}\arr\text{Dirty}
\]
然后系统在稍后：
\[
\text{Dirty}\xrightarrow{\text{writeback}}\text{Storage}
\]
这使应用不必为每次小写立即等待慢设备，也允许多次更新合并。

\subsection{三个时间点必须分开}
\[
\boxed{T_{return}}
\qquad
\boxed{T_{writeback}}
\qquad
\boxed{T_{durable}}
\]
\begin{itemize}
  \item $T_{return}$：系统调用对应用返回；
  \item $T_{writeback}$：内核把 dirty 数据提交到底层；
  \item $T_{durable}$：数据达到所要求的非易失持久边界。
\end{itemize}
存储设备自身还可能存在 volatile write-back cache，因此“设备说完成”也不必然等于“已落到非易失介质”。

\begin{warn}[\code{write()} 成功返回不等于持久化完成]
这是文件系统、页缓存 和存储设备层次最重要的边界之一。应用若需要明确的持久性语义，需要使用相应 durability 接口；文件系统/块层还可能需要 flush/FUA 等机制把设备易失缓存推进到非易失介质。
\end{warn}

\subsection{\texorpdfstring{\code{fsync()}}{fsync()} 不是简单的 Write-through}
Write-through 是缓存策略术语，强调每次写缓存同时推进 backing store；\code{fsync()} 更适合作为一个\kw{显式持久性边界请求}理解：调用方要求该文件相关的数据及必要元数据达到接口规定的持久状态。两者不能机械画等号。

\subsection{回收 页缓存：问“丢掉以后怎样恢复”}
\begin{itemize}
  \item \kw{Clean file-backed page}：文件中已有权威副本，通常可直接丢；
  \item \kw{Dirty file-backed page}：先写回或保留，不能无条件丢；
  \item \kw{Anonymous page}：没有普通文件作为后备对象，回收往往涉及 swap/压缩/其他策略。
\end{itemize}
所以不要背固定优先级，而要问：
\[
\boxed{\text{If evicted, how can this content be reconstructed? What is the cost?}}
\]

\begin{bridge}[与虚拟内存专题的接口]
页缓存 与 VM 在 \code{mmap} 处直接汇合：文件页可以通过页表映射进进程虚拟地址空间。普通 \code{read()} 常见路径是 页缓存 $\to$ 用户匿名 buffer 的显式复制；file-backed \code{mmap()} 则可以让文件页直接参与虚拟地址映射。二者最终都依赖页表、缺页和物理页生命周期。
\end{bridge}

\section{第 9 章：Spooling——把独占慢设备虚拟化成可排队服务}

\subsection{Spooling：先把任务写入磁盘队列，再由后台进程使用设备}
Spooling 的核心可以写成：
\[
\boxed{
\begin{aligned}
\text{Exclusive / Slow Physical Device}
&\arr
\text{Persistent / Buffered Queue}\\
&\arr
\text{Virtual Shared Service}
\end{aligned}
}
\]
经典 SPOOLing 依赖三个条件：多道程序使输入/输出服务与用户程序并发，通道或等价卸载机制使外围传输不必由 CPU 全程搬运，高速且可随机访问的外存保存输入/输出井。井在磁盘、缓冲区在内存，两者串联而非二选一；一个井文件对应一个任务，多个井文件再组成队列，避免不同进程的数据失去任务边界。
应用把任务提交到 spool 文件/输出队列后，不必持有物理设备直到最终完成；后台 writer/daemon 根据策略逐个把任务送到真实设备。

\subsection{经典输出 Spooling 架构模型}
\begin{center}

\begingroup
\small
\begin{tikzpicture}[>=Stealth, node distance=1.6cm and 3.0cm]
\tikzset{
  s1/.style={draw=mainblue, ultra thick, fill=softblue, rounded corners=2.5mm, align=center, minimum width=2.8cm, minimum height=0.95cm, font=\small\bfseries},
  s2/.style={draw=purple, ultra thick, fill=softpurple, rounded corners=2.5mm, align=center, minimum width=2.8cm, minimum height=0.95cm, font=\small\bfseries},
  s3/.style={draw=deeporange, ultra thick, fill=softorange, rounded corners=2.5mm, align=center, minimum width=2.8cm, minimum height=0.95cm, font=\small\bfseries},
  s4/.style={draw=deepgreen, ultra thick, fill=softgreen, rounded corners=2.5mm, align=center, minimum width=2.8cm, minimum height=0.95cm, font=\small\bfseries}
}
\node[s1] (app) {用户应用进程\\[2pt]\small User Application};
\node[s2, right=of app] (spool) {磁盘 Spooling 井\\[2pt]\small Input/Output Well};
\node[s3, right=of spool] (daemon) {守护进程 Daemon\\[2pt]\small Writer / Spool Daemon};
\node[s4, below=1.8cm of daemon] (dev) {物理独占设备\\[2pt]\small Physical Device (Printer)};

\draw[->, ultra thick, deepblue] (app) -- node[above=3pt, font=\small\bfseries] {快速写入任务} (spool);
\draw[->, ultra thick, purple] (spool) -- node[above=3pt, font=\small\bfseries] {排队读取} (daemon);
\draw[->, ultra thick, deeporange] (daemon) -- node[right=4pt, font=\small\bfseries] {慢速驱动设备} (dev);
\draw[->, ultra thick, deepgreen] (app.south) -- ++(0,-0.8cm) -| node[pos=0.25, below=3pt, font=\small\bfseries] {立即解耦返回 (无需等待物理完成)} (spool.south);
\end{tikzpicture}
\endgroup
\end{center}

\begin{itemize}
  \item \kw{时间解耦}：提交任务与物理完成不要求同时；
  \item \kw{资源解耦}：多个任务可以先排队，不必同时争用物理设备。
\end{itemize}

\subsection{408 经典输入/输出井模型}
经典 SPOOLing 系统常把磁盘上的中转区称为输入井、输出井，并由输入/输出进程负责慢速外围设备与高速磁盘之间的任务搬运。考试关注的是：
\begin{itemize}
  \item 用磁盘空间模拟独占设备的共享提交能力；
  \item 用户进程与慢设备解耦；
  \item I/O 请求进入队列，可被调度；
  \item 设备仍然一次服务有限任务，物理吞吐上限不会凭空提高。
\end{itemize}

\begin{warn}[Spooling 不是“打印物理速度变快”]
它通常不会改变打印机每秒能打印多少页；它改变的是\kw{谁需要等待、任务何时提交、设备如何排队、CPU/用户进程能否继续做其他工作}。因此它提高系统并发性和资源利用率，并显著缩短提交方占用时间，但不要把“提交完成”误当成“物理任务完成”。
\end{warn}

\begin{eng}[思想迁移，而非概念偷换]
消息队列、异步任务系统等也有“Producer $\to$ Queue $\to$ Consumer”的解耦结构，与 Spooling 在设计思想上同构。但正式术语上不应把所有队列系统都直接叫 SPOOLing。
\end{eng}

\section{第 10 章：设备分配与 I/O 调度——多个请求如何分享有限设备}

\subsection{独占、共享与虚拟设备}
从资源管理角度，设备可以按使用语义理解：
\begin{itemize}
  \item \kw{独占设备}：同一时刻通常只能被一个任务占有，如经典打印机模型；
  \item \kw{共享设备}：可以通过请求队列交错服务多个进程，如磁盘；
  \item \kw{虚拟设备}：利用 Spooling 等机制让独占物理设备表现出可并发提交的逻辑共享接口。
\end{itemize}

\subsection{设备分配涉及的资源对象}
一次 I/O 可能需要的不只是“设备本体”，还可能涉及：
\[
\text{Device}\quad+\quad\text{Controller}\quad+\quad\text{Channel/Queue Resource}
\]
经典设备分配题应明确系统采用静态分配还是动态申请、是否允许排队、是否可能形成资源等待环。

\subsection{408 经典设备分配数据结构：把“设备链”画出来}
一些教材用以下表结构描述从系统到具体硬件资源的管理关系：
\begin{itemize}
  \item \kw{SDT（System Device Table）}：系统设备总表，记录设备总体信息及其 DCT 入口；
  \item \kw{DCT（Device Control Table）}：具体设备状态、等待队列、驱动入口，并关联控制器；
  \item \kw{COCT（Controller Control Table）}：控制器状态及其连接设备/通道关系；
  \item \kw{CHCT（Channel Control Table）}：I/O 通道状态和相关资源。
\end{itemize}
于是经典分配路径可以抽象为：
\[
\boxed{\text{Process Request}\arr\text{SDT}\arr\text{DCT}\arr\text{COCT}\arr\text{CHCT}}
\]
实际 OS 实现未必使用这些名字，但考试模型借此强调：\kw{申请一个设备可能隐含申请控制器、通道和队列等多层资源}。

完整动态分配应逐层检查设备、控制器与通道是否可用；任一层忙时把进程挂入对应等待队列，全部可用才设置占用者和状态。回收则沿 DCT $\to$ COCT $\to$ CHCT 释放，并在相应资源转空闲时唤醒队首。安全分配可一次取得所需层级、避免请求且保持，但降低利用率；不安全的按需分配提高并行性，却必须用安全性或死锁处理机制约束。

设备独立性通过 LUT 把逻辑设备名映射到物理设备及驱动入口；同类设备替换和 I/O 重定向因此不要求应用改名。引入 SPOOLing 后，用户获得的不是打印机本体，而是输出井盘块与请求表；物理设备只由后台服务占用，这解释了“虚拟分配”如何移除用户进程间对打印机的直接互斥争用。

\subsection{静态分配与动态分配}
\begin{itemize}
  \item \kw{静态分配}：作业运行前一次性取得所需设备，结束后统一释放；实现简单、可避免某些运行时死锁，但利用率低；
  \item \kw{动态分配}：运行中按需申请/释放，利用率高，但必须处理等待、分配策略和潜在死锁。
\end{itemize}
如果多个进程按不同顺序申请设备/控制器/通道，就应立即联想到并发专题中的\kw{资源等待图与死锁条件}。

\subsection{为什么块设备需要调度}
当多个请求同时存在：
\[
R_1,R_2,\ldots,R_n\arr\boxed{\text{One/Finite Device Queue}}
\]
就需要策略决定服务顺序。核心评价可能包括：
\begin{itemize}
  \item 平均响应/等待时间；
  \item 吞吐；
  \item 公平性与饥饿；
  \item 设备物理成本；
  \item 优先级/截止时间。
\end{itemize}

\section{第 11 章：HDD 成本模型与经典磁盘调度}

\subsection{LBA：先把块编号与设备物理几何分开}

上层块 I/O 更希望看到从 0 开始的逻辑块地址序列，而不是让每个调用者直接操作柱面、磁头和扇区：
\[
\boxed{
\text{file offset/LBN}
\arr\text{filesystem block mapping}
\arr\text{device LBA}
\arr\text{controller/firmware physical placement}
}
\]
在经典固定几何题中，可以按题设在 CHS 与 LBA 间换算；在现代设备上，控制器还可能重映射坏块、执行磨损均衡或内部调度，因此
\[
\boxed{\text{consecutive LBA}\nRightarrow\text{guaranteed consecutive physical placement}.}
\]
这条边界也解释了为什么“文件逻辑块连续”“设备逻辑地址连续”“磁盘表面物理相邻”是三个不同判断。

\subsection{先有物理成本，再有调度策略}
经典机械硬盘访问时间：
\[
\boxed{T_{access}=T_{seek}+T_{rotation}+T_{transfer}}
\]
其中随机请求常由寻道和旋转等待主导，因此调度器试图减少磁头移动和无效等待。

若转速为 $r$ 转/秒，随机扇区的平均旋转等待为半圈：
\[
\boxed{\overline{T}_{rotation}=\frac{1}{2r}}.
\]
连续读取同一磁道的一串相邻扇区通常只在串首支付一次初始旋转等待，后续扇区随盘片旋转连续到达。调度重排主要能直接改变寻道距离；旋转位置与传输数据量是否也被纳入优化，必须看题设模型。

\subsection{六种经典算法的“设计动机”}
\begingroup
\small
\renewcommand{\arraystretch}{1.25}
\setlength{\tabcolsep}{6pt}
\begin{tabularx}{\linewidth}{>{\bfseries}p{0.18\linewidth} X X}
\toprule
算法规则 & 详细决策逻辑 & 核心设计思想与性能/饥饿风险 \\
\midrule
FCFS (先来先服务) & 按请求到达的绝对先后顺序依次服务 & 完全公平且实现简单，但磁头可能发生剧烈来回摆动 \\
SSTF (最短寻道优先) & 优先选择距离当前磁头位置最近的下一个请求 & 显著降低局部寻道时间，但两端远端请求可能严重饥饿 \\
SCAN (电梯算法) & 磁头沿一个方向移动并服务，到达端点后再反向 & 解决 SSTF 饥饿问题，改善总体移动效率与响应公平性 \\
C-SCAN (循环扫描) & 磁头仅沿单向移动并服务，到达末端后直接快速返回起点 & 服务等待时间分布比 SCAN 更均匀，避免两端不公平 \\
LOOK (LOOK 调度) & 属于 SCAN 的改进，磁头只移动到该方向最后一个请求 & 避免了 SCAN 无盲目扫到物理端点产生的多余寻道 \\
C-LOOK (C-LOOK 调度) & 属于 C-SCAN 的 LOOK 化改进 & 结合单向扫描与紧凑移动范围，不扫无请求的物理端点 \\
\bottomrule
\end{tabularx}
\endgroup

判断饥饿不要看平均距离，而要问新请求能否无限插到旧请求前。SSTF 在磁头附近持续到达新请求时可能让远端请求长期得不到服务；SCAN/LOOK 的单调扫描给出有限推进边界。磁臂粘着则是局部请求流使磁头长期停留在少数柱面，属于调度进度受局部高频请求拖住的问题，不能与某个既有请求永远不被选择的饥饿概念混写。

\begin{center}
\providecolor{themebg}{HTML}{FAFAF7}
\providecolor{themefg}{HTML}{111111}
\providecolor{themegray}{HTML}{666666}
\providecolor{themecurve}{HTML}{1D63B8}
\providecolor{themealert}{HTML}{C53030}
\providecolor{themeamber}{HTML}{B86E00}
\providecolor{themegreen}{HTML}{15803D}
\begin{tikzpicture}[font=\scriptsize,color=themefg,text=themefg,>=Stealth]
\tikzset{
  trk/.style={draw=themegray,circle,minimum size=.65cm,fill=themefg!4!themebg,font=\tiny\bfseries},
  cur/.style={trk,draw=themealert,fill=themealert!15!themebg,font=\tiny\bfseries,text=themealert},
  lbl/.style={font=\tiny\bfseries,fill=themebg,inner sep=1.5pt,rounded corners=1pt},
  flow_sstf/.style={->,line width=1.1pt,themecurve},
  flow_scan/.style={->,line width=1.1pt,themeamber},
  note/.style={rounded corners=4pt,draw=themegray!70,fill=themefg!3!themebg,inner xsep=8pt,inner ysep=7pt,align=left}
}

% ── 标题与例题说明 ──
\node[anchor=west,font=\bfseries\Large] at (0,8.8) {磁盘调度算法：SSTF、SCAN（电梯）、LOOK 磁头寻道轨迹对比};
\node[anchor=west,text=themegray,text width=17.5cm,align=left] at (0,8.2)
  {例题设定：磁道范围 $0 \sim 199$，当前磁头位置 $80$。磁道请求序列：$\{18, 45, 65, 95, 130, 175\}$。磁头初始向磁道增加方向移动。};

% ══════════════════════════════════════════════════════════════
% ── Part A: SSTF (最短寻道时间优先) (y = 4.6 .. 7.5) ──
% ══════════════════════════════════════════════════════════════
\node[font=\bfseries\normalsize,text=themecurve,anchor=west] at (0,7.5) {A. SSTF (最短寻道时间优先)：贪心选择最近磁道，易导致远端磁道“饥饿”};

\draw[themegray,line width=0.8pt] (1.0,6.8) -- (16.5,6.8);
\foreach \x/\t in {1.0/0, 2.4/18, 4.5/45, 6.06/65, 7.23/80, 8.4/95, 11.12/130, 14.63/175, 16.5/199} {
  \draw[themegray] (\x,6.7) -- (\x,6.9);
  \node[font=\tiny,above=1pt,text=themegray] at (\x,6.9) {\t};
}

\node[cur] (s_80) at (7.23,5.9) {80};
\node[trk] (s_65) at (6.06,5.9) {65};
\node[trk] (s_45) at (4.5,5.9) {45};
\node[trk] (s_18) at (2.4,5.9) {18};
\node[trk] (s_95) at (8.4,5.9) {95};
\node[trk] (s_130) at (11.12,5.9) {130};
\node[trk] (s_175) at (14.63,5.9) {175};

\draw[flow_sstf] (s_80.north) to[out=135,in=45] node[above=1pt,lbl,text=themecurve] {15} (s_65.north);
\draw[flow_sstf] (s_65.north) to[out=135,in=45] node[above=1pt,lbl,text=themecurve] {20} (s_45.north);
\draw[flow_sstf] (s_45.north) to[out=135,in=45] node[above=1pt,lbl,text=themecurve] {27} (s_18.north);
\draw[flow_sstf] (s_18.south) to[out=-20,in=-160] node[pos=0.5,below=1pt,lbl,text=themecurve] {77} (s_95.south);
\draw[flow_sstf] (s_95.north) to[out=45,in=135] node[above=1pt,lbl,text=themecurve] {35} (s_130.north);
\draw[flow_sstf] (s_130.north) to[out=45,in=135] node[above=1pt,lbl,text=themecurve] {45} (s_175.north);

\node[anchor=west,font=\tiny\bfseries,text=themecurve] at (0,4.1)
  {SSTF 访问序列：$80 \to 65 \to 45 \to 18 \to 95 \to 130 \to 175$。总寻道长度 $= 15 + 20 + 27 + 77 + 35 + 45 = 219$。};

% ══════════════════════════════════════════════════════════════
% ── Part B: SCAN vs LOOK 轨迹 (y = 1.0 .. 3.6) ──
% ══════════════════════════════════════════════════════════════
\node[font=\bfseries\normalsize,text=themeamber,anchor=west] at (0,3.5) {B. SCAN（电梯算法）vs LOOK：区别在于“是否必须到达物理尽头 199”};

\draw[themegray,line width=0.8pt] (1.0,2.8) -- (16.5,2.8);
\foreach \x/\t in {1.0/0, 2.4/18, 4.5/45, 6.06/65, 7.23/80, 8.4/95, 11.12/130, 14.63/175, 16.5/199} {
  \draw[themegray] (\x,2.7) -- (\x,2.9);
  \node[font=\tiny,above=1pt,text=themegray] at (\x,2.9) {\t};
}

% SCAN 轨迹
\draw[flow_scan] (7.23,2.3) -- (16.5,2.3)
  node[pos=0.5,above=1pt,lbl,text=themeamber] {SCAN 向外扫至物理端点 199 ($80 \to 95 \to 130 \to 175 \to 199$)};
\draw[flow_scan] (16.5,1.8) -- (2.4,1.8)
  node[pos=0.5,below=1pt,lbl,text=themeamber] {折返向内扫 ($199 \to 65 \to 45 \to 18$)};

% LOOK 轨迹标注
\node[anchor=west,font=\tiny\bfseries,text=themegreen] at (0,1.1)
  {LOOK 优化：到达最外侧请求磁道 $175$ 即可直接折返，无需移动到物理端点 $199$。总寻道长度 $= (175 - 80) + (175 - 18) = 252$。};

% ── 底部总结与误读警示 ──
\node[note,anchor=west,text width=17.5cm] at (0,-0.1)
  {\textbf{四大磁盘调度算法对比精要：}\\
   1. \textbf{FCFS}：公平无饥饿，但平均寻道距离极长；\\
   2. \textbf{SSTF}：局部最优贪心，可能导致远离磁头的请求永久饥饿；\\
   3. \textbf{SCAN / LOOK}：单向扫描到底/到最后请求再折返，消除饥饿；\\
   4. \textbf{C-SCAN / C-LOOK (循环扫描)}：单向提供服务，返回时不响应请求，提供更均匀的等待时间分布。};

\node[anchor=west,font=\scriptsize\bfseries,text=themealert,text width=17.5cm,align=left] at (0,-1.2)
  {判题陷阱：题目若未明确说明到达物理边界（0 或 199），标准考研题目默认为 SCAN 到达物理端点，LOOK 到达最远请求点。};
\end{tikzpicture}

\end{center}

\begin{method}[磁盘调度题五步法]
\begin{enumerate}
  \item 在数轴上标出当前磁头和所有请求；
  \item 明确当前方向（SCAN/C-SCAN/LOOK 类尤其关键）；
  \item 写出服务顺序，而不是边算边猜；
  \item 累加相邻位置距离；
  \item 最后检查有没有把“物理端点”与“最后一个请求位置”混淆。
\end{enumerate}
\end{method}

\subsection{磁盘组织与管理：调度之外还要知道哪些系统动作}
经典 OS 还会关注磁盘从“裸介质”变成可用存储的几个步骤：
\begin{itemize}
  \item \kw{低级/物理格式化}：建立介质扇区等物理记录结构（现代设备大量由厂商完成）；
  \item \kw{分区}：把逻辑地址空间划分成若干区域；
  \item \kw{逻辑格式化}：在分区上建立文件系统元数据结构；
  \item \kw{引导块/启动信息}：为系统启动提供初始加载路径；
  \item \kw{坏块管理}：检测不可可靠使用的介质区域，并通过重映射、备用块等机制隔离。
\end{itemize}
这些内容与下一本“文件系统专题”有接口：\kw{I/O/块设备层提供可读写的逻辑块空间，文件系统再把这些块组织成 inode、目录和文件。}

\subsection{SSD 之后为什么经典电梯思想变弱}
SSD 没有机械寻道和旋转等待，所以：
\[
T_{seek}\approx 0,\qquad T_{rotation}\approx 0
\]
经典“最短磁头移动”不再是主要目标。但现实块层仍可能需要处理：
\begin{itemize}
  \item 请求合并和队列深度；
  \item 延迟、公平与优先级；
  \item 多队列并发；
  \item 设备内部并行度。
\end{itemize}

\subsection{SSD 内部机制：没有寻道，不等于没有代价}
NAND Flash 通常以 page 为读写单位、以 block 为擦除单位，并存在原地覆盖限制。FTL (Flash Translation Layer) 维护
\[
\boxed{LBA\arr Physical\ Flash\ Page}
\]
映射，使更新可以写到新页并把旧页标为无效。空间不足时，垃圾回收选择含无效页的 block，把仍有效的页搬走后整块擦除；因此主机写入量 $W_h$ 会放大为 NAND 实际写入量 $W_n$：
\[
\boxed{WA=\frac{W_n}{W_h}\ge1}.
\]
GC、映射元数据、预留空间与随机小写都会影响写放大和尾延迟。磨损均衡把擦写分散到不同 block：动态均衡分配新写入，静态均衡还会搬迁长期不变的冷数据，目标是避免少数块先耗尽寿命。SLC/MLC/TLC/QLC 每单元位数增加通常带来更高密度和更低单位成本，也常伴随更小耐久度/更复杂读写管理；这是一组工程权衡，不应压成“SSD 一定比 HDD 在所有指标上更好”。

OS 面对 SSD 时不再以磁头距离为主要排序目标，但仍要管理请求合并、队列深度、公平、discard/TRIM 与写回节奏。连续 LBA 也不保证连续物理闪存页，因为 FTL 可随时重映射。
因此 408 的 FCFS/SSTF/SCAN 是\kw{建立“策略必须服从设备成本模型”思想的经典载体}，不是现代所有 SSD 的唯一调度逻辑。

\begin{eng}[工程边界：Linux]
现代 Linux 块层支持 mq-deadline、BFQ、Kyber、none 等 I/O scheduler；选择取决于设备和目标。工程实现会变化，但问题保持不变：\kw{多个 I/O 请求在有限服务能力下如何排序，并权衡吞吐、延迟和公平性}。
\end{eng}

\section{第 12 章：综合题——一次阻塞 \texorpdfstring{\code{read()}}{read()} 的完整路径}

\subsection{场景设定}
假设进程 A 执行：
\begin{lstlisting}
read(fd, buf, 4096);
\end{lstlisting}
文件描述符合法，但所需文件页\kw{不在页缓存（页缓存）}，底层设备采用 DMA，调用采用经典阻塞式语义。

\subsection{完整流程架构图}
\begin{center}

\begingroup
\small
\begin{tikzpicture}[>=Stealth, node distance=1.6cm and 2.4cm]
\tikzset{
  m1/.style={draw=mainblue, ultra thick, fill=softblue, rounded corners=2.5mm, align=center, minimum width=2.6cm, minimum height=0.95cm, font=\small\bfseries},
  m2/.style={draw=purple, ultra thick, fill=softpurple, rounded corners=2.5mm, align=center, minimum width=2.6cm, minimum height=0.95cm, font=\small\bfseries},
  m3/.style={draw=deeporange, ultra thick, fill=softorange, rounded corners=2.5mm, align=center, minimum width=2.6cm, minimum height=0.95cm, font=\small\bfseries},
  m4/.style={draw=deepred, ultra thick, fill=softred, rounded corners=2.5mm, align=center, minimum width=2.6cm, minimum height=0.95cm, font=\small\bfseries},
  m5/.style={draw=teal, ultra thick, fill=softteal, rounded corners=2.5mm, align=center, minimum width=2.6cm, minimum height=0.95cm, font=\small\bfseries},
  m6/.style={draw=deepgreen, ultra thick, fill=softgreen, rounded corners=2.5mm, align=center, minimum width=2.6cm, minimum height=0.95cm, font=\small\bfseries}
}
\node[m1] (s1) {1. Syscall 请求陷入\\[2pt]\small A: User $\rightarrow$ Kernel};
\node[m2, right=of s1] (s2) {2. 页缓存 Miss\\[2pt]\small 构造块请求与 DMA};
\node[m3, right=of s2] (s3) {3. 进程 A 阻塞\\[2pt]\small 切换到进程 B};

\node[m4, below=of s3] (s4) {4. 存储设备 DMA\\[2pt]\small Data $\rightarrow$ 页缓存};
\node[m5, left=of s4] (s5) {5. 设备中断 / ISR 唤醒\\[2pt]\small A: Blocked $\rightarrow$ Ready};
\node[m6, left=of s5] (s6) {6. A 重新调度 Running\\[2pt]\small Copy to User \& Return};

\draw[->, ultra thick, deepblue] (s1) -- node[above=3pt, font=\small\bfseries] {发起 \code{read}} (s2);
\draw[->, ultra thick, purple] (s2) -- node[above=3pt, font=\small\bfseries] {启动 DMA 传输} (s3);
\draw[->, ultra thick, deeporange] (s3) -- node[right=4pt, font=\small\bfseries] {硬件传输} (s4);
\draw[->, ultra thick, deepred] (s4) -- node[above=3pt, font=\small\bfseries] {触发 IRQ} (s5);
\draw[->, ultra thick, teal] (s5) -- node[above=3pt, font=\small\bfseries] {调度选中 A} (s6);
\draw[->, ultra thick, deepgreen] (s6.north) -- node[right=4pt, font=\small\bfseries] {返回用户态完成} (s1.south);
\end{tikzpicture}
\endgroup
\end{center}

\subsection{三条时间线合并}
\begingroup
\small
\renewcommand{\arraystretch}{1.25}
\setlength{\tabcolsep}{6pt}
\begin{tabularx}{\linewidth}{>{\bfseries}p{0.08\linewidth} X X X}
\toprule
时刻 & Control / CPU 状态 & Data 搬运状态 & Process State 进程状态变化 \\
\midrule
$t_0$ & A 在用户态执行 \code{read()} & 尚无数据移动 & A: User Running \\
$t_1$ & syscall 陷入内核处理 & 检查文件描述符与 offset & A: Kernel Running (内核上下文) \\
$t_2$ & 页缓存 Miss，提交块请求 & 分配/构造目标缓存页与描述符 & A: 尚未能完成返回 \\
$t_3$ & 驱动启动 DMA，调度器转去 B & 设备开始执行 DMA 数据传输 & A: Blocked, B: Running \\
$t_4$ & CPU 运行进程 B / C 的其他指令 & Device $\xrightarrow{\text{DMA}}$ 页缓存 & A: Blocked (在等待队列中) \\
$t_5$ & 设备完成传输，触发中断进入 ISR & 数据已完全入主存 Cache 页 & A: Blocked $\rightarrow$ Ready (进入就绪队列) \\
$t_6$ & 未来 CPU 调度器重新选中进程 A & - & A: Ready $\rightarrow$ Running \\
$t_7$ & 内核完成路径执行 & 页缓存 $\xrightarrow{\text{copy}}$ User Buffer & A: Kernel Running \\
$t_8$ & 系统调用返回用户态 & 用户进程可直接读 buf 数据 & A: User Running \\
\bottomrule
\end{tabularx}
\endgroup

\begin{corebox}[这条路径连接的机制]
\[
\boxed{
\begin{aligned}
&\text{Syscall}+\text{Process State}+\text{Scheduler}+\text{VFS/File}+\text{页缓存}\\
&\qquad +\text{Block Layer}+\text{Driver}+\text{DMA}+\text{Interrupt}
\end{aligned}}
\]
能够独立模拟这条路径，才算能说明 Polling、DMA、Spooling 在一次 I/O 中分别由谁执行、改变什么状态。
\end{corebox}

\section{第 13 章：概念边界}

\begingroup
\footnotesize
\renewcommand{\arraystretch}{1.22}
\begin{longtable}{>{\bfseries}p{0.15\linewidth} >{\centering\arraybackslash}p{0.025\linewidth} >{\bfseries}p{0.15\linewidth} p{0.31\linewidth} p{0.25\linewidth}}
\toprule
对象 A & $\neq$ & 对象 B & 真正区别与题目信号 & 混淆后果\\
\midrule
\endfirsthead
\toprule
对象 A & $\neq$ & 对象 B & 真正区别与题目信号 & 混淆后果\\
\midrule
\endhead
System Call & $\neq$ & Interrupt & 前者通常由程序主动请求内核服务；后者通常由外部异步硬件事件触发。两者都可能进入内核，但事件来源不同。 & 把“进入内核”误当成同一种事件，进而错判同步性、返回点与处理路径。\\
Interrupt & $\neq$ & Context Switch & 中断只给内核一次处理机会；是否切换执行实体还取决于后续状态变化与调度决策。 & 把每次中断都计为一次进程切换。\\
Polling & $\neq$ & PIO & Polling 说明怎样发现设备状态变化；PIO 说明 CPU 是否亲自执行设备数据搬运。两者可以组合。 & 把控制发现方式和数据传输方式混成一个维度。\\
Interrupt & $\neq$ & Async I/O & 设备可以用中断报告完成，但用户接口仍可能同步阻塞；异步性描述调用者是否需要等待结果。 & 看到中断就误判用户调用不会阻塞。\\
DMA & $\neq$ & Zero-copy & DMA 卸载设备与内存之间的一段搬运；zero-copy 还要求减少其他内存复制和数据重排。 & 低估用户/内核复制、IOMMU 与 cache 同步等剩余成本。\\
Buffer & $\neq$ & Cache & Buffer 主要吸收速率或粒度不匹配；Cache 主要保留未来可能复用的慢对象副本。 & 不能解释数据取走后是否仍值得保留，以及命中/替换策略为何存在。\\
Hardware Cache & $\neq$ & 页缓存 & 前者属于硬件存储层级；后者由 OS 按文件对象与页索引管理文件内容副本。 & 把地址翻译/Cache 行为与文件缓存状态混在同一 Owner。\\
SPOOLing & $\neq$ & Buffer & SPOOLing 是中转存储、请求队列与后台服务的组合，用于把独占设备虚拟成可排队服务。 & 把设备虚拟化问题降格成“多放一块缓冲区”。\\
\code{write()} 返回 & $\neq$ & 持久化完成 & buffered write 返回通常只表示内核已接受数据；持久化还涉及 writeback、提交与设备稳定介质。 & 把接口返回时间错误当成断电后数据必然存在的时间。\\
\code{fsync()} & $\neq$ & Write-through & 前者是应用显式请求的持久边界接口；后者是缓存写策略。 & 用缓存策略术语替代具体系统调用的完成承诺。\\
I/O Complete & $\neq$ & Process Running & 完成通常使等待条件成立并令任务 Blocked $\rightarrow$ Ready；获得 CPU 仍由 scheduler 决定。 & 把 wakeup 直接画成 Running。\\
SCAN & $\neq$ & LOOK & SCAN 可继续走到物理端点；LOOK 只走到该方向最后一个待服务请求。 & 磁头移动距离与服务终点计算错误。\\
C-SCAN & $\neq$ & C-LOOK & 两者都单向服务；C-LOOK 同样只走到有请求的边界，不扫无请求的物理端点。 & 把两种循环策略的回绕距离机械算成相同。\\
\bottomrule
\end{longtable}
\endgroup

\section{第 14 章：最小调用接口}

遇到 I/O 问题时，先区分三条机制线：\kw{控制路径}回答谁提交、谁发现完成；\kw{数据路径}回答谁搬数据、经过哪些缓冲或缓存；\kw{任务状态路径}回答调用者何时继续、何时阻塞以及完成后进入什么状态。只要题目开始要求固定十三问、题型路由、时间轴计算或考场检查，转入同目录训练文档 \texttt{IO请求与完成.md}。

若问题继续追踪完整 blocking \code{read()}，转入 OS-I01；若追踪 DMA/IRQ 的硬件责任边界，转入对应跨科 Bridge。本册不在 Canonical 中重复维护训练清单。

\section{第 16 章：专题接口}

\subsection{与“进程、线程、调度与控制权”专题}
I/O 把抽象的进程状态图变成真实事件：
\[
\boxed{\text{read waits}\arr Running\to Blocked}
\]
\[
\boxed{\text{I/O completion}\arr Blocked\to Ready}
\]
\[
\boxed{\text{scheduler}\arr Ready\to Running}
\]

\subsection{与“CPU × OS 软硬件协作”专题}
\[
\boxed{
\text{Hardware: IRQ/DMA/MMIO}
\qquad
\text{Kernel: Driver/WaitQueue/Scheduler/Policy}
}
\]
I/O 是软硬件边界最密集的一块。

\subsection{与“内存虚拟化”专题}
页缓存、DMA buffer、用户 buffer、\code{mmap} 都会牵涉：
\[
\text{VA}\arr\text{PTE}\arr\text{Physical Page}
\]
但“文件内容缓存”与“地址翻译缓存”要严格分层。

\subsection{与下一本“文件系统”专题}
文件系统解决：
\[
\boxed{\text{Name/FD}\arr\text{File Object}\arr\text{Logical File Block}}
\]
本册解决：
\[
\boxed{\text{I/O Request}\arr\text{Transfer}\arr\text{Completion}}
\]
两者在 页缓存、块层和 \code{read/write/mmap} 路径处汇合。

\section{核心结论：I/O 是异步状态转换系统}
传统章节容易留下这样的知识碎片：
\begin{quote}
Polling、Interrupt、DMA、Buffer、Spooling、Disk Scheduling 各背一段。
\end{quote}
本册希望建立的是另一种模型：
\[
\boxed{
\begin{aligned}
\text{谁提出请求？}&\arr\text{谁接管控制？}\arr\text{谁搬数据？}\\
&\arr\text{谁等待？}\arr\text{完成怎样被发现？}\arr\text{状态怎样恢复？}
\end{aligned}}
\]
进一步再问：
\[
\boxed{
\begin{aligned}
&\text{能否缓存？能否批量？能否排队？能否卸载？}\\
&\text{瓶颈在 CPU、内存带宽、设备、队列还是持久性边界？}
\end{aligned}}
\]
只要能沿这两组问题完整推演，I/O 就不再是操作系统里最零散的一章，而会变成连接\kw{进程状态、调度、软硬件边界、虚拟内存、文件系统与存储设备}的中心枢纽。

\newpage
\section*{参考与工程边界资料}
\addcontentsline{toc}{section}{参考与工程边界资料}
本册核心结构服务于 408 教学模型；以下资料用于验证现代实现边界，避免把教材抽象误认为唯一工程实现。
\begin{enumerate}
  \item Linux Kernel Documentation, \emph{Entry/exit handling for exceptions, interrupts, syscalls and KVM}.\\
  \url{https://docs.kernel.org/core-api/entry.html}
  \item Linux Kernel Documentation, \emph{Dynamic DMA mapping Guide / DMA API}.\\
  \url{https://docs.kernel.org/core-api/dma-api-howto.html}\\
  \url{https://docs.kernel.org/core-api/dma-api.html}
  \item Linux Kernel Documentation, \emph{Buffered I/O / iomap operations}.\\
  \url{https://docs.kernel.org/filesystems/iomap/operations.html}
  \item Linux Kernel Documentation, \emph{Explicit volatile write back cache control}.\\
  \url{https://docs.kernel.org/block/writeback_cache_control.html}
  \item Linux Kernel Documentation, \emph{Block layer / I/O schedulers}.\\
  \url{https://docs.kernel.org/block/}
  \item Linux Kernel Documentation, \emph{Driver model / driver API}.\\
  \url{https://docs.kernel.org/driver-api/driver-model/overview.html}
  \item IBM Documentation, \emph{Output spooling / output queues}.\\
  \url{https://www.ibm.com/docs/en/i/7.5?topic=concepts-output-spooling}
  \item Remzi H. Arpaci-Dusseau and Andrea C. Arpaci-Dusseau, \emph{Operating Systems: Three Easy Pieces}, I/O Devices chapters.\\
  \url{https://pages.cs.wisc.edu/~remzi/OSTEP/}
\end{enumerate}

\end{document}
